\scanpage{35}

and $H^0$ are left exact, one deduces from this a homomorphism of exact diagrams of sets
% No arrow strokes are visible between these six objects.
\diagram{
\begin{tikzcd}[ampersand replacement=\&,column sep=2.0em,row sep=1.7em]
H^0(X,F) \& H^{0}(X,G) \& H^0(X,H) \\
H^0(Y,F_Y) \& H^0(Y,G_Y) \& H^0(Y,H_Y),
\end{tikzcd}
}
where according to what precedes the vertical arrows are injective\inserted{(in particular the last}, and where by hypothesis the middle vertical arrow is bijective. (NB the subscripts $Y$ denoting the restrictions to $Y$)? One formally concludes from this that the first vertical arrow is also bijective.
% The inserted parenthesis runs into the right edge; no closing parenthesis is visible.

We shall apply this remark to the case \deleted{where}
\[
G=h\ h\ (F)\ x\ (\quad(F))\ ,
\]
% Spaces in this and the following formulas reproduce unfilled places in the typescript.
and to the obvious homomorphism $F\qquad G$, deduced from the natural homomorphisms into the two factors. This homomorphism is injective, as follows immediately (by consideration of the fibers) from the hypothesis $X=Y\quad(\overline{X})$. It remains to prove the assertion of bijectivity for $G$, or what amounts to the same, for the two factors, hence for sheaves of the form $h(\unsure{H|\mathbb{H}})$ (resp. $(\overline{F})$),where $H$ (resp. $\overline{F}$) is a sheaf on $Y$ (resp. on $\overline{X}$). now the assertion is trivial for a sheaf of the form $h(H)$. For the case of a sheaf of the form $(\overline{F})$, taking into account the fact that the base c\deleted{y}hange homomorphism $H(\quad(\overline{F}))\qquad{}_{Y}(\overline{h}(\overline{F}))$ is bijective, one sees that the homomorphism $H^0(X,\quad(\overline{F}))\qquad H^0(Y,h\quad(\overline{F}))$ is identified with the homomorphism $H^0(\overline{X},\overline{F})\qquad H^0(\overline{Y},\overline{h}\ (\overline{F}))$, which is bijective by hypothesis. This completes the proof of (i).

(ii) The conclu\deleted{\unsure{t|b}}sion o\deleted{'}f\deleted{inject} bijectivity for $H^0(X,F)\qquad H^0(Y,h\ (F))$ \margin{\handwritten{$\unsure{\backslash|\checkmark}$}}for every sheaf of sets $F$ on $X$, already obtained in (i), ensures \deleted{by an immediate argument} that for every sheaf of groups $F$ on $X$ (not necessarily ind-finite), th\deleted{'}e\deleted{homomorphism of the} restriction functor
% archive sheet number: 34
